% F085｜真实数学图；代表函数不构成抽象解的唯一性证明。
\begin{center}
\begin{tikzpicture}
\begin{axis}[width=12.4cm,height=5.1cm,axis lines=middle,
 xmin=-2.3,xmax=3.5,ymin=-7,ymax=10.5,
 xtick={-2,-1,0,1,2,3},xticklabels=\empty,ytick={-6,-3,0,3,6,9},xlabel={$t$},ylabel={$f(t)$},
 title={正比例直线 $f(t)=3t$},title style={font=\small,yshift=3pt},
 tick label style={font=\small},label style={font=\small},
 ylabel style={at={(axis description cs:0,1)},anchor=south east,xshift=-18pt},
 clip=false,grid=major,grid style={black!8},enlargelimits=false,
 every axis plot/.append style={thick},samples=161]
\addplot[CExplanation,very thick,domain=-2:3.1] {3*x};
\draw[dashed,black!35] (axis cs:1,0) |- (axis cs:0,3);
\draw[dashed,black!35] (axis cs:2,0) |- (axis cs:0,6);
\draw[dashed,black!35] (axis cs:3,0) |- (axis cs:0,9);
\draw[dashed,black!35] (axis cs:-1,0) |- (axis cs:0,-3);
\addplot[only marks,mark=*,mark size=2pt,orange!80!black] coordinates {(1,3) (2,6) (3,9) (-1,-3)};

\node[font=\small,anchor=north,yshift=-3pt] at (axis cs:-2,-7) {$-2$};
\node[font=\small,anchor=north,yshift=-3pt] at (axis cs:-1,-7) {$-1$};
\node[font=\small,anchor=north,yshift=-3pt] at (axis cs:1,-7) {$1$};
\node[font=\small,anchor=north,yshift=-3pt] at (axis cs:2,-7) {$2$};
\node[font=\small,anchor=north,yshift=-3pt] at (axis cs:3,-7) {$3$};

\end{axis}
\end{tikzpicture}
\par\small 输入相加，高度相加；点 $(1,3)$ 与 $(-1,-3)$ 关于原点对称。
\end{center}
